\section{Problem formulation as an MDP}
\label{sec:mdp}

After every drafted token the agent chooses \CONT{} or \STOP{}. Its reward is the number of tokens gained minus the time cost of the round, measured in tokens.

\begin{table}[h]
\centering
\caption{MDP definition.}
\label{tab:mdp}
\begin{tabularx}{\linewidth}{@{}lX@{}}
\toprule
Part & Definition \\
\midrule
Episode & Generating the full answer (e.g.\ 128 tokens) for one prompt, made of many draft$\to$verify rounds \\
State & $s$ = tokens drafted so far this round (1--8); $p$ = top probability of the latest draft token; $H$ = entropy of the draft distribution; $q$ = product of $p$ over this round (estimated chance that all drafts survive); $a$ = moving average of the recent acceptance rate \\
Actions & \CONT{} (draft one more token) or \STOP{} (send to the target). \STOP{} is forced at $s = K = 8$ \\
Reward & 0 on \CONT{}. On \STOP{}: $r = (\min(s,n)+1) - \lambda(1 + s\,c)$ \\
Discount & $\gamma = 1$ (episodic and bounded), or $0.99$ for stability with DQN and REINFORCE \\
\bottomrule
\end{tabularx}
\end{table}

Here $c = t_\text{draft} / t_\text{target}$ is the measured cost of one draft step relative to one target call, and $\lambda$ is the price of one target call in tokens. \todo{measured values of $c$ for each model pair.}

\paragraph{Choosing $\lambda$.} Maximising throughput (tokens per unit time) is a ratio objective. By Dinkelbach's method~\citep{dinkelbach1967} it equals maximising $\text{tokens} - \lambda\cdot\text{time}$, where $\lambda$ is the best achievable throughput. In practice we set $\lambda$ to the throughput of the best fixed-$k$ baseline, train, measure the new throughput, update $\lambda$ and retrain once. We also report a small sweep over $\lambda$.

\paragraph{Bellman optimality equations.}
\begin{align}
Q^*(s,\STOP) &= \mathbb{E}\left[\min(s,n)+1\right] - \lambda(1+s\,c) + \gamma\,\mathbb{E}\left[V^*(s_{\text{next round}})\right] \\
Q^*(s,\CONT) &= \gamma\,\mathbb{E}\left[\max_{a'} Q^*(s',a')\right] \\
V^*(s) &= \max\left(Q^*(s,\STOP),\; Q^*(s,\CONT)\right)
\end{align}
For the tabular methods the next-round term is the same under both actions, so each round can be treated as a short episode of at most 8 steps.
